\documentclass[10pt,a4paper]{amsart}
\usepackage[margin=19mm]{geometry}
\usepackage{amsmath,amssymb,mathtools}
\usepackage[scr=rsfs,cal=euler]{mathalfa}
\usepackage{xfrac}
\usepackage[unicode,hidelinks,bookmarks=false]{hyperref}
\newcommand{\ie}{i.e.\ }
\newcommand{\cf}{cf.\ }
\newcommand{\resp}{respectively\ }
\newcommand{\Subsection}{\subsection}
\providecommand{\nobreakdash}{\nobreak\hbox{-}}
\setlength{\emergencystretch}{2em}
\multlinegap0pt
\usepackage{tikz}
\usetikzlibrary{cd}
\usepackage{amssymb}
\usepackage[shortlabels]{enumitem}
\usepackage[normalem]{ulem}
\usepackage{xcolor}
\newtheorem{setup}{Setup}
\newtheorem{lemma}{Lemma}
\usepackage{cleveref}
\crefname{setup}{Setup}{Setups}
\crefname{lemma}{Lemma}{Lemmas}
\title{Fiberwise criteria for Fourier--Mukai equivalences}
\author{El\'ias Guisado Villalgordo, Pat Lank, Kabeer Manali Rahul, Nebojsa Pavic}
\date{Source: arXiv:2512.16503v4 (CC BY 4.0)}
\begin{document}
\maketitle

\noindent\emph{Source and licence.} Fiberwise criteria for Fourier--Mukai equivalences. Version arXiv:2512.16503v4 (CC BY 4.0), distributed by arXiv under the Creative Commons Attribution 4.0 International licence, http://creativecommons.org/licenses/by/4.0/, as recorded in the arXiv licence metadata for this exact version and shown on the abstract page https://arxiv.org/abs/2512.16503v4. Full licence notice: \texttt{licenses/arxiv-2512.16503-CC-BY-4.0.txt}.\par\smallskip

\noindent\textit{Editorial scope.} Counted unit: the article's lemma lem:fm\_cube\_pullback\_natural\_isomorphism (source0.tex lines 1103--1110). The article's complete original proof of it is delivered with it (source0.tex lines 1112--1152, one \texttt{proof} environment of 41 lines). No mathematics is written, repaired or re-derived: the statement and the proof are the article's own lines, in the article's own notation, and no retained line is substituted or re-flowed.  Retained uncounted as the source-local context the proof needs: lines 622--648.  Nothing was omitted from any retained span, so the package carries no omission record.  No retained line contains a citation, so the wrapper carries no bibliography and no bibliography entry.  The retained lines contain the article's own diagram code, which the wrapper typesets with the standard tikz packages; no image file is delivered and the package declares no asset dependency.  Editorial formatting only, in addition: an AMS \texttt{amsart} preamble that restates the article's own declarations and macros used by the retained lines, namely the environments \texttt{setup}, \texttt{lemma} on separate counters, together with the standard packages those lines need.  The retained environments print in this document as Setup 1, Lemma 1 in order of appearance, whereas the article prints the same items with the numbering of its own full document.  The complete native source of the article as distributed in the arXiv e-print for this exact version is bundled unchanged as \texttt{sources/191242/source0.tex}.
\begin{setup}
    \label{setup:fm_cube_pullback}
    Let $t\colon T \to S$ be a morphism of Noetherian schemes. Suppose $f_i\colon Y_i \to S$ are proper flat morphisms with $i\in \{1,2\}$. Consider the commutative diagram
    \begin{equation}
        \label{diag:fm_cube_pullback}
        % https://q.uiver.app/#q=WzAsOCxbNCwxLCJUIl0sWzQsMywiUyJdLFsxLDMsIllfMSJdLFsyLDIsIllfMiJdLFsxLDEsIllfMVxcdGltZXNfUyBUIl0sWzIsMCwiWV8yXFx0aW1lc19TIFQiXSxbMCwwLCJZXzFcXHRpbWVzX1MgWV8yXFx0aW1lc19TIFQiXSxbMCwyLCJZXzFcXHRpbWVzX1MgWV8yIl0sWzAsMSwidCJdLFsyLDEsImZfMSIsMl0sWzMsMSwiZl8yIl0sWzQsMCwiZ18xIl0sWzUsMCwiZ18yIl0sWzQsMiwidF8xIiwwLHsibGFiZWxfcG9zaXRpb24iOjcwfV0sWzUsMywidF8yIiwyLHsibGFiZWxfcG9zaXRpb24iOjgwLCJzdHlsZSI6eyJib2R5Ijp7Im5hbWUiOiJkYXNoZWQifX19XSxbNiw0LCJnXzJeXFxwcmltZSIsMl0sWzYsNSwiZ18xXlxccHJpbWUiLDJdLFs3LDMsImZeXFxwcmltZV8xIiwwLHsibGFiZWxfcG9zaXRpb24iOjMwLCJzdHlsZSI6eyJib2R5Ijp7Im5hbWUiOiJkYXNoZWQifX19XSxbNywyLCJmXlxccHJpbWVfMiIsMl0sWzYsNywidF5cXHByaW1lIiwyXV0=
        \begin{tikzcd}
            {Y_1\times_S Y_2\times_S T} && {Y_2\times_S T} \\
            & {Y_1\times_S T} &&& T \\
            {Y_1\times_S Y_2} && {Y_2} \\
            & {Y_1} &&& S
            \arrow["{g_1^\prime}"', from=1-1, to=1-3]
            \arrow["{g_2^\prime}"', from=1-1, to=2-2]
            \arrow["{t^\prime}"', from=1-1, to=3-1]
            \arrow["{g_2}", from=1-3, to=2-5]
            \arrow["{t_2}"'{pos=0.8}, dashed, from=1-3, to=3-3]
            \arrow["{g_1}", from=2-2, to=2-5]
            \arrow["{t_1}"{pos=0.7}, from=2-2, to=4-2]
            \arrow["t", from=2-5, to=4-5]
            \arrow["{f^\prime_1}"{pos=0.3}, dashed, from=3-1, to=3-3]
            \arrow["{f^\prime_2}"', from=3-1, to=4-2]
            \arrow["{f_2}", from=3-3, to=4-5]
            \arrow["{f_1}"', from=4-2, to=4-5]
        \end{tikzcd}
    \end{equation}
    obtained by base change along $t$.
\end{setup}

\begin{lemma}
	\label{lem:fm_cube_pullback_natural_isomorphism}
	Consider \Cref{setup:fm_cube_pullback}. Let $K\in D_{\operatorname{qc}}(Y_1\times_S Y_2)$. Then on $D_{\operatorname{qc}}$ there is a natural isomorphism
    \begin{displaymath}
        \alpha^K \colon \Phi_K \circ \mathbf{R}(t_1)_\ast 
        \to \mathbf{R}(t_2)_\ast \circ \Phi_{\mathbf{L}(t^\prime)^\ast K}.
    \end{displaymath}
\end{lemma}

\begin{proof}
	This follows from the string of natural isomorphisms for each $E\in D_{\operatorname{qc}}(Y_1\times_S T)$:
    \begin{align*}
        \Phi_K \circ \mathbf{R} (t_1)_\ast (E)
        &= \mathbf{R} (f_1^\prime)_\ast 
        \big( \mathbf{L} (f^\prime_2)^\ast \mathbf{R} (t_1)_\ast E \otimes^{\mathbf{L}} K \big)
        \\
        &\cong 
        \mathbf{R} (f_1^\prime)_\ast 
        \big( 
        \mathbf{R} t^\prime_\ast \mathbf{L}(g^\prime_2)^\ast E 
        \otimes^{\mathbf{L}} K 
        \big)
        &
        \text{(flat base change)}
        \\
        &\cong 
        \mathbf{R} (f_1^\prime)_\ast 
        \mathbf{R} t^\prime_\ast 
        \big( 
        \mathbf{L}(g^\prime_2)^\ast E 
        \otimes^{\mathbf{L}} \mathbf{L}(t^\prime)^\ast K 
        \big) 
        &
        \text{(projection formula)}
        \\
        &\cong 
        \mathbf{R} (t_2)_\ast \mathbf{R} (g^\prime_1)_\ast
        \big( 
        \mathbf{L}(g^\prime_2)^\ast E 
        \otimes^{\mathbf{L}} \mathbf{L}(t^\prime)^\ast K 
        \big)
        &
        \text{(pseudofunctoriality of $\mathbf{R}(-)_\ast$)}
        \\
        &= 
        \mathbf{R} (t_2)_\ast 
        \circ 
        \Phi_{\mathbf{L}(t^\prime)^\ast K} (E). && \qedhere
    \end{align*}
\end{proof}

\end{document}
